\section{\bi{From collision stability to exact contacts}{由碰撞穩定性到精確接觸}}
\begin{thm}\label{th:lock}
\bi{Let $N=2m\ge4$, $C=1+3\pi m^2$, $R=\sum_{n\ge2} n^4q^{n^2-1}\le2/3$. If $0<\eta\le1/(NC^2)$, all minimizers consist of two packed progressions of steps $h$, antipodal midpoints and counts $2k+1,2l+1$, where $k+l=m-1$, $|k-l|\le1$. Each outer gap is $1/2-(m-1)h$. There is one species-preserving isometry orbit for $k=l$, two for $k\ne l$, exchanged by global sign reversal. The width condition holds for $q\le17/50$. For $N=2$ under the same width condition the antipodal pair minimizes uniquely up to rotation for every feasible floor.}{設 $N=2m\ge4$、$C=1+3\pi m^2$、$R=\sum_{n\ge2}n^4q^{n^2-1}\le2/3$。若 $0<\eta\le1/(NC^2)$，全部極小者由兩列公差恰為 $h$ 的緊簇組成，中點對徑，粒子數分別為 $2k+1,2l+1$，其中 $k+l=m-1$、$|k-l|\le1$。兩個外間隙均為 $1/2-(m-1)h$。在保號等距變換下，$k=l$ 時有一個軌道，$k\ne l$ 時有兩個，整體換號互換二者。$q\le17/50$ 足以滿足寬度條件。對 $N=2$，在相同寬度條件下，任一可行下限均由對徑對唯一取等，模整體平移。}
\end{thm}
\begin{proof}
\bi{Write $a_n=q^{n^2-1}$ and $K(u)=\sum_{n\ge1}a_n\cos(nu)$. The ratio bound $R\le16q^3/(1-(3/2)^4q^5)<2/3$ holds at and below $17/50$. Under the stated $R$ hypothesis, $\|K''\|_\infty,\|K^{(4)}\|_\infty<2$, $K''(\pi)\ge1-R/4\ge5/6$, and $D=K''(\pi)-K''(0)=2\sum_{n\ {\rm odd}}n^2a_n$ obeys $2\le D<4$. All sums over parity below use positive frequencies.}{置 $a_n=q^{n^2-1}$、$K(u)=\sum_{n\ge1}a_n\cos(nu)$。相鄰項比值給出 $R\le16q^3/(1-(3/2)^4q^5)<2/3$，適用於 $q\le17/50$。在定理的 $R$ 條件下，$\|K''\|_\infty,\|K^{(4)}\|_\infty<2$、$K''(\pi)\ge1-R/4\ge5/6$、$D=K''(\pi)-K''(0)=2\sum_{n\ {\rm odd}}n^2a_n$ 滿足 $2\le D<4$。奇偶頻率和均只取正整數。}

\bi{\emph{Localization.} Let $A$ consist of the arcs from a positive site to the next negative site, so $\mu=(\mathbf1_A)'=\sum\sigma_j\delta_{\theta_j}$. Integration by parts twice gives $Q_v(A)=\iint v(u-w)\dd\mu(u)\dd\mu(w)=\int_{A^2}-v''(u-w)\dd u\dd w$. For $v=K-\cos/3$, both $v$ and $-v''$ decrease weakly on $[0,\pi]$: the first-mode derivative dominates the tails by $|\sin(nu)|\le n\sin u$. Add a constant to $-v''$ to make it positive. Circle rearrangement \cite[Thm.~4.1]{circle} bounds the integral by that on an arc of equal length $b$; the constant cancels. Its value is $2[v(0)-v(b)]$, maximal at $b=\pi$. The reserved kernel contributes $|z|^2/3$, $z=\int_A e^{iu}\dd u$. With $H$ a semicircle and $I_N=1-2[G_s(0)-G_s(1/2)]/N$, this yields}{\emph{定位。} 令 $A$ 為從正粒子至下一負粒子的有向弧之聯集，則 $\mu=(\mathbf1_A)'=\sum\sigma_j\delta_{\theta_j}$。兩次分部積分得 $Q_v(A)=\iint v(u-w)\dd\mu(u)\dd\mu(w)=\int_{A^2}-v''(u-w)\dd u\dd w$。取 $v=K-\cos/3$，由 $|\sin(nu)|\le n\sin u$ 及尾項界，$v$ 和 $-v''$ 在 $[0,\pi]$ 上皆非增。為 $-v''$ 加常數使其為正，圓周重排 \cite[定理 4.1]{circle} 給出同長度 $b$ 的單弧上界，並可抵消所加常數。單弧的值 $2[v(0)-v(b)]$ 在 $b=\pi$ 最大。保留的一階核貢獻 $|z|^2/3$，其中 $z=\int_A e^{iu}\dd u$。記半圓為 $H$，$I_N=1-2[G_s(0)-G_s(1/2)]/N$，則}
\begin{equation}\label{eq:deficit}
 F=\frac{N(E_s-I_N)}{2cq}=Q_K(H)-Q_K(A)\ge (4-|z|^2)/3 .
\end{equation}
\bi{Choose $H$ where $\cos(u-\arg z)>0$ (any phase for $z=0$). Then $J=\int_{A\triangle H}|\cos(u-\arg z)|\dd u=2-|z|\le3F/2$. Placing measure $b$ nearest the two zeros gives $J\ge4[1-\cos(b/4)]$ for $b=|A\triangle H|$. Thus at zero floor equality requires and suffices that the signed jumps reduce to an antipodal dipole. Splitting an odd cluster about one end gives strictly ordered approximants, so this is the open-class infimum. For $N>2$ equality needs collisions, whereas the two-site dipole is feasible at all floors.}{取 $\cos(u-\arg z)>0$ 所在的半圓 $H$（$z=0$ 時相位任取）。則 $J=\int_{A\triangle H}|\cos(u-\arg z)|\dd u=2-|z|\le3F/2$。把長度 $b$ 集中在兩零點附近可得 $J\ge4[1-\cos(b/4)]$，其中 $b=|A\triangle H|$。故在零下限的閉包中，帶號跳躍抵消後恰為對徑偶極是取等的充要條件。將一端的奇數簇拆開即有嚴格有序的逼近列，所以 $I_N$ 也是開類下確界。$N>2$ 時取等須有碰撞，兩粒子對徑構型則在每個可行下限取等。}

\bi{For a packed odd block put $B_j(w)=\sum_{p=-j}^j(-1)^{j+p}\cos(pw)$. Pairing consecutive cosines gives $|1-B_j(w)|\le j(j+1)w^2/2$, since $2|\cos((r-1)w)-\cos(rw)|\le(2r-1)w^2$. The antipodal balanced trial has squared amplitude $[B_k(n\eta)-(-1)^n B_l(n\eta)]^2$. Even modes cannot increase its deficit and $4-y^2\le4(2-y)$ for odd-mode amplitude $y=B_k+B_l$. Hence $F_{\rm trial}\le D A_m\eta^2$, $A_m=k(k+1)+l(l+1)=\lfloor m^2/2\rfloor$. Any minimizer satisfies $J\le3m^2\eta^2$. Rotate $H$ to $(0,\pi)$. If a jump lies distance $r\ge\eta$ from the nearer end, one of the two adjacent intervals of length $\eta/2$ belongs to the discrepancy. Since $\sin u\ge2u/\pi$ for $0\le u\le\pi/2$, $J\ge(\eta/\pi)(r-\eta/2)$. Every site thus lies within $C\eta$ of an endpoint. Outside the two neighborhoods the indicators agree: otherwise an intervening middle quarter-circle contributes $\sqrt2>3m^2\eta^2$. Here $C\eta<\pi/4$. The net jumps in the two neighborhoods are $+1,-1$, so both blocks are odd.}{令緊奇數簇的振幅 $B_j(w)=\sum_{p=-j}^j(-1)^{j+p}\cos(pw)$。相鄰餘弦配對，利用 $2|\cos((r-1)w)-\cos(rw)|\le(2r-1)w^2$，得 $|1-B_j(w)|\le j(j+1)w^2/2$。均衡對徑試探的振幅平方為 $[B_k(n\eta)-(-1)^n B_l(n\eta)]^2$。偶模不會增加虧損，奇模取 $y=B_k+B_l$ 有 $4-y^2\le4(2-y)$。故 $F_{\rm trial}\le D A_m\eta^2$，$A_m=k(k+1)+l(l+1)=\lfloor m^2/2\rfloor$。每個極小者有 $J\le3m^2\eta^2$。將 $H$ 旋至 $(0,\pi)$。若跳點距較近端點 $r\ge\eta$，則其左右長 $\eta/2$ 的兩段中有一段全屬差異集；在 $[0,\pi/2]$ 用 $\sin u\ge2u/\pi$ 得 $J\ge(\eta/\pi)(r-\eta/2)$。故各粒子距端點至多 $C\eta$。兩鄰域外指示函數一致，否則中間四分之一圓貢獻 $\sqrt2>3m^2\eta^2$；此處 $C\eta<\pi/4$。兩端鄰域的淨跳躍是 $+1,-1$，各簇均為奇數。}

\bi{\emph{Selection and contact.} In each cluster use offsets $u_i$ from its assigned antipode and normalize the initial sign to $+1$. Set $P=\sum(-1)^iu_i$, $V=\sum(-1)^iu_i^2-P^2$. Write the gaps $a_1,b_1,\ldots,a_j,b_j$. Consecutive differences of squares and cancellation of $P^2$ give}{\emph{分簇與接觸。} 以各自的對徑端點為原點，記偏移 $u_i$，將各簇的起始符號歸一為正。置 $P=\sum(-1)^iu_i$、$V=\sum(-1)^iu_i^2-P^2$，簇內間隙依次為 $a_1,b_1,\ldots,a_j,b_j$。相鄰平方作差並減去 $P^2$ 得}
\begin{equation}\label{eq:mom}
 P=u_0+\sum_r b_r,\qquad V=2\sum_{r=1}^j b_r\sum_{i=1}^r a_i\ge j(j+1)\eta^2 .
\end{equation}
\bi{Each gap derivative of $V$ is at least $2\eta$. In the Taylor expansion of each ordered pair at $0$ or $\pi$, odd derivatives vanish and the offset is at most $2C\eta$. The signed sum of squared differences within a cluster is $2V$; across the two clusters with initial signs normalized it is $V_0+V_\pi+(P_0-P_\pi)^2$. Counting both cross directions gives}{每個間隙的 $V$ 偏導數皆至少為 $2\eta$。在 $0$ 或 $\pi$ 處對各有向粒子對作 Taylor 展開，奇階導數為零，偏移不逾 $2C\eta$。同簇的帶號差平方和為 $2V$，異簇起始符號各自歸一後為 $V_0+V_\pi+(P_0-P_\pi)^2$。將異簇兩個方向計入，得}
\begin{equation}\label{eq:taylor}
 F=D(V_0+V_\pi)+K''(\pi)(P_0-P_\pi)^2+\mathcal R,\qquad
 |\mathcal R|\le\tfrac43 N^2C^4\eta^4 .
\end{equation}
\bi{An unbalanced pair increases the sum of the two $j(j+1)$ by at least two. The error is at most $4\eta^2/3<2D\eta^2$, contradicting the trial bound. To vary gap $p$ at fixed signed first moments, use velocity $\mathbf1_{i\ge p}-\mathbf1_{p\ {\rm even}}$ in that cluster, zero outside. A small decrease of a slack gap draws its suffix or prefix inward, preserving the other internal gaps and widening one outer gap. All relative velocities have modulus at most one. Expanding $K'$ with remainder $\|K^{(4)}\|_\infty(2C\eta)^3/6$ per pair gives $\partial_p F\ge2D\eta-(8/3)N^2 C^3\eta^3>0$. A slack gap contradicts minimality.}{若兩簇不均衡，兩項 $j(j+1)$ 的和至少增加二。餘項至多 $4\eta^2/3<2D\eta^2$，與試探上界矛盾。改變第 $p$ 個間隙但固定兩簇帶號一階矩，取本簇粒子速度 $\mathbf1_{i\ge p}-\mathbf1_{p\ {\rm even}}$，其餘為零。將鬆弛間隙稍微縮短只令其前綴或後綴內移，保持其他內間隙並增大一個外間隙。相對速度絕對值至多為一；直接展開 $K'$，逐對餘項至多 $\|K^{(4)}\|_\infty(2C\eta)^3/6$，故 $\partial_p F\ge2D\eta-(8/3)N^2 C^3\eta^3>0$。鬆弛間隙與極小性矛盾。}

\bi{\emph{Midpoints.} In a packed odd block $P$ equals its midpoint offset. At separation $\pi+v$ the variable part of $F$ is twice the alternating cross sum of $K(\pi+v+(j-i)\eta)$, even in $v$ by reflection. On the segment to zero every argument stays within $2C\eta$ of $\pi$ by convexity, and both outer gaps stay feasible. The signed coefficients have sum one before doubling. Bounding the change in each $K''$ by $\|K^{(4)}\|_\infty(2C\eta)^2/2$ and the number of terms by $N^2/4$ gives $F_{vv}\ge2K''(\pi)-2/C^2>0$. Hence $v=0$. Equal counts allow a half turn to interchange signs; unequal counts distinguish the size carrying positive net charge, even under reflection. Existence and sign symmetry realize both assignments.}{\emph{中點。} 在緊奇數簇中 $P$ 等於中點偏移。當中點相隔 $\pi+v$，$F$ 的可變部分為 $K(\pi+v+(j-i)\eta)$ 之交替跨簇和的兩倍。反射使其為 $v$ 的偶函數。沿 $v$ 到零的線段，各核自變量距 $\pi$ 不逾 $2C\eta$，兩個外間隙保持可行，兩者均由端點及凸性得到。倍乘前各係數的帶號和為一。每項 $K''$ 的變化至多 $\|K^{(4)}\|_\infty(2C\eta)^2/2$，項數不逾 $N^2/4$，故 $F_{vv}\ge2K''(\pi)-2/C^2>0$，必有 $v=0$。兩簇同大時平移半圈互換電荷；不同大時，帶正淨電荷的簇大小在保號反射下仍可區分。極小者的存在性與換號對稱性保證兩種指派皆取等。}
\end{proof}

\begin{cor}\label{th:energy}
\bi{Under Theorem~\ref{th:lock}, put $A=k(k+1)$, $B=l(l+1)$ and $S_{p,o},S_{p,e}=\sum n^p q^{n^2}$ over odd and even $n$, respectively. Then}{在定理~\ref{th:lock} 的條件下，記 $A=k(k+1)$、$B=l(l+1)$，並以 $S_{p,o},S_{p,e}$ 分別表示 $\sum n^p q^{n^2}$ 的奇、偶頻率和。則}
\begin{align}
 M_N &=1-\frac{2c}{N}\sum_{n\ge1}q^{n^2}[B_k(n\eta)-(-1)^nB_l(n\eta)]^2,\label{eq:exact}\\
 M_N &= I_N+\frac{4c(A+B)}N S_{2,o}\eta^2 \notag\\
 &\quad-\frac{2c}{N}\left[\left(\frac{(A+B)^2}{4}+
 \frac{A(A-1)+B(B-1)}6\right)S_{4,o}+\frac{(A-B)^2}4 S_{4,e}\right]\eta^4+O(\eta^6).\label{eq:fourth}
\end{align}
\end{cor}
\begin{proof}
\bi{Insert the minimizing coordinates in \eqref{eq:energy}. Pairing the outer terms in the finite exponential sum gives $B_j(u)=2\cos(ju)-B_{j-1}(u)$, $B_0=1$. Induction gives $B_k(u)=1-Au^2/2+A(A-1)u^4/24+O(u^6)$. Squaring the sum and difference yields \eqref{eq:fourth}. The trigonometric polynomials have bounded sixth derivatives, and the sixth Gaussian moment is finite at fixed $s,N$. For six sites $A=B=2$, hence the quadratic and quartic coefficients are $(8c/3)S_{2,o}$ and $-(14c/9)S_{4,o}$.}{將極小構型代入式 \eqref{eq:energy}。取有限指數和的兩端得遞推 $B_j(u)=2\cos(ju)-B_{j-1}(u)$、$B_0=1$。歸納得 $B_k(u)=1-Au^2/2+A(A-1)u^4/24+O(u^6)$，分別對和、差平方即得式 \eqref{eq:fourth}。三角多項式的六階導數有界，固定 $s,N$ 時高斯六階矩有限。六粒子有 $A=B=2$，二次、四次係數分別為 $(8c/3)S_{2,o}$ 與 $-(14c/9)S_{4,o}$。}
\end{proof}

